% introduction.tex -- Batch Aggregation Paper (BNR-2026-004)

\section{Introduction}

Enterprise blockchain deployments face a fundamental tension between data privacy
and verifiable correctness. Organizations in regulated industries require that
operational data---supply chain events, industrial process records, financial
transactions---remain confidential, yet auditors and counterparties must be able
to verify the integrity of state transitions without accessing the underlying data.
Zero-knowledge validium architectures address this tension by maintaining enterprise
state off-chain while submitting succinct validity proofs to a shared settlement layer,
achieving privacy through data withholding and correctness through cryptographic
guarantees~\cite{chaliasos2024analyzing}.

A critical but underexplored component of validium node architecture is the
\emph{batch aggregation layer}: the subsystem responsible for collecting individual
transactions, persisting them durably, and forming fixed-size batches suitable for
ZK circuit proving. Unlike public rollup sequencers that can tolerate transaction
reordering or soft-confirmation semantics, enterprise validium nodes must provide
\emph{zero transaction loss} under arbitrary crash scenarios---a single lost
industrial traceability record can invalidate an entire audit chain. Furthermore,
batch formation must be deterministic (identical transaction sets always produce
identical batches) to support replay-based verification and regulatory compliance.

This paper makes three contributions. First, we present a WAL-based (Write-Ahead Log)
batch aggregation protocol that combines size-based and time-based batch formation
triggers in a HYBRID strategy, achieving bounded latency under low load and bounded
batch size under high load. Second, we provide comprehensive benchmarks demonstrating
274,000~tx/min peak throughput with sub-millisecond batch formation latency, 450/450
determinism tests passing, and 150/150 crash recovery tests with zero transaction loss.
Third---and most significantly---we report how TLA+ model checking of our protocol
discovered a \emph{silent transaction loss bug} that 150 empirical crash recovery
tests failed to detect. The bug manifests during a narrow window (1.9--12.8 seconds)
between batch formation and ZK proof completion; a crash in this window causes
irrecoverable data loss with no error indication. We present the formal diagnosis,
the corrected protocol, and its exhaustive verification.

The key finding is methodological: formal verification and empirical testing are
complementary, not redundant. TLC model checking explored the full state space
(2,630 distinct states) in under one second and found a 6-step counterexample that
no reasonable test suite would have covered---not because the test suite was
insufficient, but because the failure requires reasoning about the \emph{absence}
of an action (crash at a specific interleaving point) rather than the presence of
observable behavior.
